\documentclass[11pt,a4paper]{article}

\usepackage[utf8]{inputenc}
\usepackage[T1]{fontenc}
\usepackage{amsmath,amssymb}
\usepackage{geometry}
\usepackage{graphicx}
\usepackage{hyperref}
\usepackage{booktabs}
\usepackage{enumitem}

\geometry{margin=2.5cm}

\title{\textbf{Nonlinear Stabilisation of the Electron Soliton}\\[4pt]
	\large Documentation for the OpenWave M4 Implementation}
\author{Łukasz Smoliński}
\date{\today \\ \small  Version: 1.0.00}
	
\begin{document}
	\maketitle
	
	\begin{abstract}
		This document provides a self-contained technical specification for the
		nonlinear wave equation that is proposed to stabilise the electron core
		($K_{WC}=10$) in the vector M4 model of the OpenWave platform.  The
		coupling coefficient is derived from the BCC lattice geometry of Energy
		Wave Theory, three complementary forms of the nonlinear term are presented,
		and the role of the Degraded EMC Wall as a necessary boundary condition is
		quantified.  The 1-3-6 wave-centre arrangement is shown to be the unique
		self-consistent candidate satisfying the topological and energetic
		constraints of the vacuum lattice.  The stability of this configuration
		is a postulate whose numerical verification is the immediate objective
		of the M4 implementation described herein.
	\end{abstract}
	
	%=====================================================================
	\section{Context and Objective}
	%=====================================================================
	

	
The scalar M3 model of OpenWave has demonstrated that spherical
phyllotaxis (golden-angle, $\sim 137.5^{\circ}$) resolves the
K-selectivity problem at the tested neighbouring values: $K=10$ forms
a stable standing wave (amplitude flat at $\sim 0.29$ after $\sim 140$
steps under a $0.1$ perturbation), while $K=9$ and $K=11$ decay
systematically ($-34\%$ and $-45\%$, respectively, with $K=9$
confirmed across two independent runs).\footnote{Pull Request \#205,
	\url{https://github.com/openwave-labs/openwave/pull/205}, merged;
	closes/contributes to Issue \#203.}  In the vector M4 model, however,
the same golden-angle arrangement does \emph{not} stabilise the
electron core, indicating that the missing ingredient is not purely
geometric but nonlinear.
	
	
	%=====================================================================
	\section{General Soliton Wave Equation}
	%=====================================================================
	
	The scalar amplitude $\Psi(\mathbf{r},t)$ of the standing wave is
	governed by
	\begin{equation}\label{eq:soliton}
		\Bigl( \frac{\partial^{2}}{\partial t^{2}} - c^{2}\nabla^{2} \Bigr)
		\Psi(\mathbf{r},t) + \mathcal{F}(\Psi) = 0 ,
	\end{equation}
	where $c$ is the speed of light and $\mathcal{F}$ encodes the nonlinear
	self-interaction that prevents dispersive collapse.
	
	%=====================================================================
	\section{Coupling Coefficient from BCC Geometry}
	%=====================================================================
	
	In EWT the magnetic deficit $\epsilon_{M}$ is the fundamental lattice
	response parameter.  It is defined through the calibrated stiffness
	$N_{\text{final}}$ and the geometric limit $8\pi^{4}$:
	\begin{equation}\label{eq:epsM}
		\epsilon_{M} = \frac{1}{N_{\text{final}}\,\pi^{3}}
		\;\approx\; \frac{1}{8\pi^{7}}, \qquad
		N_{\text{final}} = 778.818123 .
	\end{equation}
	The approximate equality reflects the $0.058\%$ lattice impedance
	$\delta$ between $N_{\text{final}}$ and the ideal BCC packing limit
	$8\pi^{4}$.  Physically, $\epsilon_{M}$ measures the fractional loss of
	stiffness caused by the soliton's spin-induced magnetic torque.
	
	The natural coupling strength for the nonlinear term is the inverse
	deficit,
	\begin{equation}\label{eq:gamma}
		\gamma \equiv \frac{1}{\epsilon_{M}}
		= N_{\text{final}}\,\pi^{3} \approx 2.4148 \times 10^{4}.
	\end{equation}
	For numerical implementation two variants are available:
	\begin{itemize}
		\item $\gamma_{\text{geo}} = 8\pi^{7} \approx 2.4162 \times 10^{4}$
		(ideal BCC limit);
		\item $\gamma_{\text{final}} = N_{\text{final}}\pi^{3} \approx 2.4148
		\times 10^{4}$ (anchored to the measured $\alpha$).
	\end{itemize}
	The $0.058\%$ difference may be critical for discriminating $K=10$ from
	neighbouring configurations.  The M4 implementation should initially use
	$\gamma_{\text{final}}$ and, if necessary, explore the geometric limit
	as a consistency check.
	
	%=====================================================================
	\section{Three Complementary Forms of $\mathcal{F}$}
	%=====================================================================
	
	\subsection*{Variant A --- Pure cubic nonlinearity}
	The simplest form compatible with NLS-type soliton equations is
	\begin{equation}\label{eq:F_cubic}
		\boxed{\mathcal{F}(\Psi) = \gamma\,\Psi^{3}} .
	\end{equation}
	A term proportional to $\Psi^{3}$ is the lowest-order odd nonlinearity
	that preserves the sign of the restoring force; it arises from a quartic
	potential $V(\Psi)=\tfrac{1}{4}\gamma\Psi^{4}$.
	
	\subsection*{Variant B --- Density-modulated nonlinearity}
	Let $\rho(r)$ be the local EMC density inside the soliton and $\rho_{0}$
	the statutory vacuum background.  The deficit fraction is
	$(1-\rho(r)/\rho_{0})$, and the nonlinear term becomes
	\begin{equation}\label{eq:F_density}
		\boxed{\mathcal{F}(\Psi,\rho)
			= \gamma\,\Bigl(1-\frac{\rho(r)}{\rho_{0}}\Bigr)\,\Psi^{3}} .
	\end{equation}
	The nonlinearity is maximal in the core where $\rho(r) \ll \rho_{0}$ and
	vanishes in the surrounding vacuum, linking stability directly to the
	EMC push-out mechanism.
	
	\subsection*{Variant C --- Explicit 1-3-6 source-driven dynamics}
	In the native EWT model the electron consists of ten discrete, mobile
	wave centres (WCs) arranged in a 1-3-6 configuration.  The vector wave
	equation then takes the driven form
	\begin{equation}\label{eq:F_source}
		\boxed{\Bigl( \frac{\partial^{2}}{\partial t^{2}} - c^{2}\nabla^{2} \Bigr)
			\mathbf{\Psi} + \gamma\,\|\mathbf{\Psi}\|^{2}\,\mathbf{\Psi}
			= \mathbf{J}_{1\text{-}3\text{-}6}(\mathbf{r},t)} ,
	\end{equation}
	where the source operator is the sum of ten vectorial centres:
	\begin{equation}\label{eq:J}
		\mathbf{J}_{1\text{-}3\text{-}6}(\mathbf{r},t) =
		\mathbf{j}_{\text{core}}(\mathbf{r} - \mathbf{r}_{0})
		+ \sum_{i=1}^{3} \mathbf{j}_{\text{inner}}(\mathbf{r} - \mathbf{r}_{i}(t))
		+ \sum_{j=1}^{6} \mathbf{j}_{\text{outer}}(\mathbf{r} - \mathbf{r}_{j}(t)) .
	\end{equation}
	Each class carries a distinct phase offset; the phase asymmetry between
	the inner (3) and outer (6) groups forces a synchronised rotation,
	generating the $720^{\circ}$ spin of the electron.
	
	%=====================================================================
	\section{Structural Emergence of the 1-3-6 Geometry}
	%=====================================================================
	
	Why the 1-3-6 partition?  Three complementary reasons are proposed
	(pending numerical verification):
	
	\begin{enumerate}
		\item The \textbf{core} (1) anchors the soliton at a single BCC lattice
		node, providing the central phase reference.
		\item The \textbf{inner shell} (3) forms an equilateral triangle ---
		the minimal number of vertices that can define a plane whose
		normal coincides with the spin axis.  It is proposed that four
		or more centres on a single shell at this radius would introduce
		phase frustration.
		\item The \textbf{outer shell} (6) is proposed to complete the
		octahedral coordination of the BCC unit cell, saturating the
		remaining propagation axes and guaranteeing isotropic far-field
		emission after spherical masking by the Degraded EMC Wall.
	\end{enumerate}
	
	It is proposed that configurations such as 2-3-5 or 2-2-6 would
	misalign the spin axis or leave BCC axes unmatched, producing a net
	multipole moment that the nonlinearity cannot cancel.  The 1-3-6
	triad is therefore the unique self-consistent candidate satisfying:
	(a) the topological winding number $Q=10$,
	(b) the cubic nonlinearity with coefficient $\gamma$, and
	(c) the octahedral symmetry of the $Im\bar{3}m$ BCC lattice.
	A rigorous proof that this configuration minimises the energy
	functional remains a task for the M4 simulation.
	
	%=====================================================================
	\section{Topological Selection Rule}
	%=====================================================================
	
	On the spherical shell $S^{2}$ that encloses the soliton, the field
	configuration is characterised by an integer winding number
	\begin{equation}\label{eq:Q}
		Q = \frac{1}{4\pi} \int_{S^{2}} \Psi^{*}\,\omega \in \mathbb{Z},
	\end{equation}
	where $\omega$ is the normalised area form on $S^{2}$.  For the electron
	$Q = K_{WC} = 10$.
	
	This number is a topological invariant --- it cannot change under
	continuous deformations of the wave field.  The configuration of
	$Q=10$ wave centres is a natural candidate for the lowest-energy
	stable homotopy class that can establish a fully isotropic boundary
	condition on $S^{2}$ (analogous to the global minimum of the spherical
	Thomson problem).  A rigorous proof that $Q=10$ yields a deeper
	potential well than $Q=8,9,11$ requires an explicit evaluation of the
	energy functional with $\gamma\Psi^{3}$ and the $Im\bar{3}m$ projection
	--- the immediate objective of the M4 simulation.
	
	Equation~\eqref{eq:soliton} together with any of the three nonlinear
	variants and $Q=10$ constitutes a well-posed, falsifiable model.  The
	stability of the 1-3-6 configuration is at present a postulate, not a
	proof.
	
	\begin{enumerate}
		\item The \textbf{nonlinearity} ($\gamma$) creates a deep, narrow
		potential well that resists dispersion.
		\item The \textbf{topological winding number} is proposed to fix the
		ground state at exactly ten wave centres, pending numerical
		confirmation that $Q=10$ is energetically preferred over
		neighbouring integers.
		\item The \textbf{BCC stiffness constant} $N_{\text{final}}$ (or its
		geometric limit $8\pi^{4}$) fixes the coupling without adjustable
		parameters.
	\end{enumerate}
	
	%=====================================================================
	\section{Dimensional Weight Operator}
	%=====================================================================
	
	Each factor of $\pi$ in the coupling budget corresponds to a distinct
	geometric degree of freedom of the soliton.  Labelling these as
	$\pi^{(1)}$ (wave amplitude), $\pi^{(2)}$ (frequency),
	$\pi^{(3)},\pi^{(4)},\pi^{(5)}$ (3D spatial substrate),
	$\pi^{(6)}$ (soliton radius), and $\pi^{(7)}$ (charge), the
	coefficient \eqref{eq:gamma} can be factorised as
	\begin{equation}\label{eq:factor}
		\gamma = \frac{1}{\epsilon_{M}} = N_{\text{final}}\,\pi^{3}
		\;\approx\; 8\pi^{7}
		= 8\;\cdot\;\pi^{(1)}\pi^{(2)}\pi^{(3)}\pi^{(4)}\pi^{(5)}\pi^{(6)}\pi^{(7)} .
	\end{equation}
	The approximate equality with $8\pi^{7}$ reflects the geometric limit
	$\gamma_{\text{geo}}$; the exact factorisation in terms of labelled
	$\pi^{(k)}$ is a formal bookkeeping device valid for any numerical
	prefactor.  The factor~8 corresponds to the eight primary propagation
	axes of the $Im\bar{3}m$ unit cell, ensuring isotropic nonlinear
	restoring force.
	
	%=====================================================================
	\section{Degraded EMC Wall}
	%=====================================================================
	
	The nonlinear stabilisation requires a boundary condition: the
	Degraded EMC Wall, where the EMC density transitions from the deep
	core deficit to the statutory background $\rho_{0}$.
	
	A natural length scale for the wall radius is provided by the
	electron Compton wavelength $\lambda_C$.  Using the standard QED
	relations $r_e = \alpha\lambda_C/2\pi$ and $1/\alpha \approx 137$,
	\begin{equation}\label{eq:rwall}
		\boxed{r_{\text{wall}} \sim \frac{\lambda_C}{2\pi}
			= \frac{r_e}{\alpha} \approx 137\,r_e} .
	\end{equation}
	Because $\alpha$ is itself derivable from BCC geometry, this scale
	introduces no new parameter.  The wall provides a broad transition
	zone ($\sim 137\,r_e$) within which the nonlinearity can adiabatically
	relax the EMC density, suppressing sharp gradients that would
	otherwise destabilise the standing wave.
	
	\subsection*{Possible EMC Density Peak}
	
	Two density profiles are compatible with the integral deficit: a
	monotonic approach to $\rho_{0}$, or a local peak
	($\rho_{\text{wall}} > \rho_{0}$) caused by EMC accumulation.  If a
	peak exists, it would assist stabilisation in two ways:
	\begin{enumerate}
		\item In Variant~B, the factor $(1-\rho/\rho_{0})$ changes sign
		within the peak, creating a local defocusing barrier that
		prevents the standing wave from leaking outward.
		\item In all variants, a density peak implies a locally higher
		elastic modulus, acting as a partial reflector that increases
		the $Q$-factor of the soliton resonator.
	\end{enumerate}
	Whether the peak is indispensable or auxiliary can only be decided
	once the full radial EMC profile is computed from the nonlinear wave
	equation in the M4 simulation.
	
	%=====================================================================
	\section{Internal Wave-Centre Dynamics (Yee Hypothesis)}
	%=====================================================================
	
	An additional hypothesis, put forward in the foundational EWT
	literature, proposes that the electron's ten wave centres are not
	static but undergo continuous micro-motion around their nodal
	positions.  A displaced centre experiences a restoring force toward
	the local amplitude minimum; this perpetual shifting is the origin of
	the $720^{\circ}$ spin and, simultaneously, a dynamic stabilisation
	mechanism.  The 1-3-6 geometry ensures that only a fraction of
	centres are off-node at any instant, while the remainder anchor the
	structure.  This hypothesis should be tested in M4 once the nonlinear
	term and the EMC Wall are in place: comparing simulations with and
	without WC motion will reveal whether internal dynamics is essential
	or secondary.
	
	%=====================================================================
	\section{Relation Between Energetic and Topological Arguments}
	%=====================================================================
	
	Two complementary conditions single out $K_{WC}=10$:
	
	\begin{enumerate}
		\item \textbf{Energetic:} The $r^{5}/r^{3}$ scaling disparity creates
		a deep potential well for intermediate wave-centre counts.  $K_{WC}$
		cannot be too small (shallow well) or too large (volumetric
		overload).
		\item \textbf{Topological:} The winding number $Q$ must be an integer.
		Only discrete values are permitted.
	\end{enumerate}
	
	These are not independent post-hoc justifications; they are mutually
	reinforcing.  The specific value $K=10$ emerges as the unique integer
	that simultaneously satisfies both constraints under the octahedral
	symmetry of the $Im\bar{3}m$ BCC lattice.  A rigorous proof that no
	other integer passes both filters remains a task for the M4 simulation.
	
	%=====================================================================
	\section{Summary and Implementation Recommendations}
	%=====================================================================
	
	\begin{itemize}
		\item \textbf{Wave equation:}
		$(\partial_{t}^{2} - c^{2}\nabla^{2})\Psi + \mathcal{F}(\Psi) = 0$.
		\item \textbf{Coupling:}
		$\gamma = 1/\epsilon_{M} = N_{\text{final}}\pi^{3} \approx 2.41\times 10^{4}$.
		Use $\gamma_{\text{final}}$ initially; the $0.058\%$ difference
		from $\gamma_{\text{geo}} = 8\pi^{7}$ may be decisive.
		\item \textbf{Three variants for $\mathcal{F}$:}
		\begin{itemize}
			\item[A.] $\gamma\Psi^{3}$ (pure cubic, minimal implementation).
			\item[B.] $\gamma(1-\rho/\rho_{0})\Psi^{3}$ (density-modulated,
			links stability to the EMC push-out deficit).
			\item[C.] $\gamma\|\mathbf{\Psi}\|^{2}\mathbf{\Psi}$ with an
			explicit 1-3-6 source operator (directly encodes the
			electron's wave-centre architecture).
		\end{itemize}
		\item \textbf{Degraded EMC Wall:}
		$r_{\text{wall}} \sim r_e/\alpha \approx 137\,r_e$.  A possible
		density peak would further assist stabilisation.
		\item \textbf{Topological selection:}
		$Q=10$ on $S^{2}$ is proposed as the selection rule; rigorous
		proof that $Q=10$ is energetically preferred remains an M4 task.
		\item \textbf{1-3-6 partition:}
		The unique self-consistent candidate; its stability is a postulate
		pending numerical verification.
		\item \textbf{Energetic + topological arguments:}
		Complementary constraints that together single out $K_{WC}=10$.
		\item \textbf{Model status:}
		Well-posed and falsifiable; stability of the 1-3-6 configuration
		is a postulate whose numerical verification is the immediate
		objective of the OpenWave M4 simulation.
	\end{itemize}
	
	\bigskip
	\begin{center}
		\rule{0.5\textwidth}{0.4pt}
		\textit{Document prepared for the OpenWave development team.}\\[4pt]
	\end{center}
	
	%=====================================================================
	\begin{thebibliography}{9}
		
		\bibitem{Smolinski2025}
		Smoliński, Ł. (2025).
		\textit{The Geometric Identity of Gravity and Dimensional Unification
			Resolving $\alpha$, Lepton $(g-2)_l$, Weinberg, and Cabibbo Mixing}.
		DOI: \href{https://doi.org/10.5281/zenodo.17654657}{10.5281/zenodo.17654657}.
		
	\end{thebibliography}
	
\end{document}